\documentclass[11pt]{article}
\usepackage[a4paper,margin=1in]{geometry}
\usepackage{amsmath,amssymb,mathtools,bm}
\usepackage{enumitem,booktabs,array}
\usepackage{tcolorbox}
\usepackage{hyperref}
\usepackage[protrusion=true,expansion=false]{microtype}
\hypersetup{colorlinks=true,linkcolor=blue,urlcolor=blue,citecolor=blue}
\setlist[itemize]{topsep=2pt,itemsep=2pt}
\setlist[enumerate]{topsep=2pt,itemsep=2pt}
\newtcolorbox{keybox}{colback=blue!3!white,colframe=blue!40!black,boxrule=0.4pt,arc=1mm}
\newtcolorbox{alertbox}{colback=red!3!white,colframe=red!40!black,boxrule=0.4pt,arc=1mm}
\newtcolorbox{answerbox}{colback=green!3!white,colframe=green!40!black,boxrule=0.4pt,arc=1mm}
\newtcolorbox{drillbox}{colback=yellow!5!white,colframe=orange!60!black,boxrule=0.4pt,arc=1mm}
\newcommand{\EV}{\mathbb{E}}
\newcommand{\Var}{\mathrm{Var}}
\newcommand{\Cov}{\mathrm{Cov}}
\newcommand{\blank}[1]{\underline{\hspace{#1}}}

\title{W05-03: Schmidt, Timmermann \& Wermers (2016, AER) \\
\large ``Runs on Money Market Mutual Funds'' --- Active Recall Workbook}
\author{Banking \& Risk Management --- Exam Prep}
\date{\today}
\begin{document}\maketitle

\begin{keybox}
\textbf{Paper in one sentence.} STW use daily share-class-level MMMF flow data during September 2008 to show that (i) \emph{institutional} MMMF share classes ran sharply, retail classes did not, (ii) funds with higher exposure to Lehman-like risky assets lost more flows, and (iii) within the same fund, institutional redeem faster, consistent with a first-mover-advantage run dynamic that maps directly to Diamond--Dybvig.

\textbf{Placement.} The canonical evidence on the 2008 MMMF run, motivating the 2010 and 2014 SEC reforms (floating NAV for institutional prime MMMFs). Complements Gorton--Metrick (W04-01) on the repo-side of the shadow-banking run.
\end{keybox}

\section*{Round 1: Complete Walk-Through}

\subsection*{1.1 The Reserve Primary event}
16 Sept 2008: Reserve Primary Fund ``breaks the buck'' after writing down Lehman holdings. Prime MMMFs hold short-term commercial paper; depositors are daily redemption at \$1.00 NAV. The run spreads across prime funds the next week, prompting a Treasury guarantee.

\subsection*{1.2 Data}
Daily share-class level flows for US prime MMMFs, roughly 400 funds. Share classes separately identify institutional from retail investors within the same fund. Holdings data from N-Q and N-MFP.

\subsection*{1.3 Identification of run behaviour}
\begin{itemize}
\item \textbf{Within-fund, institutional vs.\ retail.} Fund-day fixed effects absorb fund-specific asset risk; the difference in flow between institutional and retail share classes identifies differential run propensity.
\item \textbf{Cross-fund.} Funds with higher ex-ante exposure to Lehman or similar-risk paper lose more flows.
\end{itemize}

\subsection*{1.4 Headline results}
\begin{itemize}
\item Institutional share classes experienced $\sim$15\%+ outflows in the week after Reserve Primary broke the buck; retail classes little or no outflow.
\item Flows are concave in risk exposure --- a ``threshold'' effect consistent with global-games.
\item Large, fast institutional redemption is profitable under \$1.00 NAV: those who redeem first get par; late redeemers bear losses. This is a pecuniary \emph{first-mover advantage}.
\end{itemize}

\subsection*{1.5 Policy implications}
\begin{itemize}
\item Floating NAV for institutional prime funds (SEC 2014) directly attacks the first-mover advantage.
\item Redemption gates and liquidity fees act as ex-post tools.
\item STW quantify that a floating-NAV institutional design eliminates most of the 2008 run.
\end{itemize}

\subsection*{1.6 Caveats}
\begin{alertbox}
\begin{itemize}
\item Short event window --- external validity to future events is assumed.
\item Within-fund comparison requires investors in different share classes to hold the same underlying assets; STW argue this is enforced by fund structure.
\item The Treasury guarantee may have dampened retail outflows endogenously.
\end{itemize}
\end{alertbox}

\section*{Round 2: Level 1 Fill-in}
\begin{enumerate}
\item Trigger event: \blank{2cm} Fund breaks the buck, \blank{1cm} Sep 2008.
\item Key comparison: \blank{1.5cm} vs.\ \blank{1cm} share classes within the same fund.
\item Institutional outflow magnitude: $\sim$\blank{1cm}\%$+$ in the week.
\item First-mover advantage mechanism: fixed \blank{1cm} NAV.
\item Policy response: \blank{2cm} NAV for institutional prime funds.
\end{enumerate}

\section*{Round 3: Level 2 Prompts}
\begin{enumerate}
\item Map the MMMF run to Diamond--Dybvig. Who is the depositor and what is the illiquid asset?
\item Why does a \$1.00 NAV create a first-mover advantage? Derive it mechanically.
\item Why are retail investors less sensitive? Give two candidate reasons.
\item Discuss how floating NAV kills the run advantage.
\end{enumerate}

\section*{Round 4: Minimal Scaffolding}
From a blank page: (i) event, (ii) data, (iii) within-fund comparison, (iv) first-mover logic, (v) policy.

\section*{Round 5: Blank Page}
Essay: MMMFs as modern banks --- the 2008 run and the floating-NAV fix.

\vspace{6pt}
\begin{drillbox}
\textbf{Speed Drill.} (1) Trigger fund. (2) Within-fund comparison. (3) Institutional outflow fact. (4) First-mover-advantage source. (5) Main SEC reform.
\end{drillbox}

\section*{Quick Reference \& Answer Key}
\begin{answerbox}
\begin{itemize}
\item \textbf{Event.} Reserve Primary breaks the buck, 16 Sept 2008.
\item \textbf{Run.} Institutional share classes run; retail do not.
\item \textbf{Mechanism.} Fixed \$1.00 NAV $\Rightarrow$ first-mover advantage.
\item \textbf{Policy.} Floating NAV for institutional prime MMMFs (SEC 2014).
\end{itemize}
\end{answerbox}

\begin{keybox}
\textbf{30-second exam pitch.} Schmidt, Timmermann \& Wermers (2016) document the September 2008 run on US prime money-market mutual funds at daily share-class frequency. When Reserve Primary Fund broke the buck after writing down Lehman holdings, institutional share classes in prime MMMFs experienced rapid, large outflows, while retail share classes in the same funds did not. Funds with more exposure to Lehman-like risk lost more flows. The central mechanism is the fixed \$1.00 NAV: early redeemers get par, late redeemers absorb the losses --- a pecuniary first-mover advantage that is Diamond--Dybvig's strategic complementarity in modern form. Institutional investors, being larger and faster, exploit this first-mover advantage. The paper motivates the SEC's 2014 move to floating NAV for institutional prime MMMFs, which removes the advantage and dampens run risk.
\end{keybox}

\end{document}
